\subsection{\texorpdfstring{NORMALISER OF \(W_{a}\)}{NORMALISER OF W\_a}}

In this no., we assume that the chosen scalar product on V is invariant not only under W but under the whole of the group A(R). We identify A(R) and A(R \(^{\prime}\) ).

Let G be the normaliser of \(W_{a}\) in the group of displacements of the affine euclidean space E. If g is a displacement of E, and s is the orthogonal reflection with respect to a hyperplane L, the displacement \(gsg^{-1}\) is the orthogonal reflection with respect to the hyperplane \(g(L)\) . It follows that G is the set of displacements of E that permute the hyperplanes \(L_{\alpha,k}\) (for \(\alpha \in R\) and \(k \in Z\) ).

Now, the group of automorphisms of E is the semi-direct product of the orthogonal group U of \(V^{*}\) and the group T of translations. If \(u \in U\) and \(v \in V^{*}\) , the hyperplane \(L_{\alpha,k}\) is transformed by \(g = u \circ t(v)\) into the hyperplane with equation

\[
\langle^ {t} u ^ {- 1} (\alpha), x \rangle = k + \langle \alpha , v \rangle .
\]

Consequently, \(g \in \mathbf{G}\) if and only if, on the one hand \({}^t u\) permutes the roots, in other words belongs to \(\mathrm{A(R)}\) , and on the other hand \(\langle \alpha, v \rangle \in \mathbf{Z}\) for all \(\alpha \in \mathbb{R}\) , that is \(v \in \mathrm{P(R^{\prime})}\) . In other words, the group \(G\) is the semi-direct product of \(A(R)\) by \(P\) . Since \(Q \subseteq P\) and \(W \subseteq A(R)\) , the quotient group \(G / W_a\) is the semi-direct product of \(A(R)/W\) by \(P(R^{\prime})/Q(R^{\prime})\) ; it is immediately checked that the corresponding action of \(A(R)/W\) on \(P(R^{\prime})/Q(R^{\prime})\) is the canonical action (§1, no. 9).

We denote by \(W_{a}^{\prime}\) the subgroup of G formed by the semi-direct product of W by P. This is a normal subgroup of G, and \(G/W_{a}^{\prime}\) is canonically isomorphic to \(A(R)/W\) ; moreover, the canonical map from \(P(R^{\prime})\) to \(W_{a}^{\prime}/W_{a}\) gives by passing to the quotient an isomorphism from \(P(R^{\prime})/Q(R^{\prime})\) to \(W_{a}^{\prime}/W_{a}\) .

Now let C be an alcove of E, and let \(G_{C}\) be the subgroup consisting of the elements \(g \in G\) such that \(g(\mathrm{C}) = \mathrm{C}\) . Since \(W_{a}\) is simply-transitive on the alcoves, the group G is the semi-direct product of \(G_{C}\) by \(W_{a}\) . The corresponding isomorphism from \(G/W_{a}\) to \(G_{C}\) gives rise in particular to a canonical isomorphism from \(\mathrm{P}(\mathrm{R}^{\prime})/\mathrm{Q}(\mathrm{R}^{\prime})\) to the group \(\Gamma_{C} = G_{C} \cap W_{a}^{\prime}\) .

Assume that R is irreducible, and retain the notation of Prop. 5 of no. 2. Put \(R_{0} = R\) , and let \(R_{i} (i \in I)\) be the root system generated by the \(a_{j}\) , for \(j \neq i\) . For i = 0 (resp. \(i \in I\) ), let \(w_{i}\) be the unique element of \(W(R_{i})\) (identified with a subgroup of W) which transforms the positive roots of \(R_{i}\) relative to the basis \((\alpha_{j})_{j \neq i}\) into negative roots (§ 1, no. 6, Cor. 3 of Prop. 17).

PROPOSITION 6. For all \(i \in \mathbf{J}\) , the element \(\gamma_i = t(\overline{\omega}_i)w_i w_0\) belongs to \(\Gamma_{C}\) and the map \(i \mapsto \gamma_i\) is a bijection from \(\mathbf{J}\) to \(\Gamma_{C} - \{1\}\) .

We remark first of all that the root \(w_{i}(\tilde{\alpha})\) is of the form

\[
n _ {i} \alpha_ {i} + \sum_ {j \neq i} b _ {i j} \alpha_ {j},
\]

and hence is positive.

We show that, if \(i \in \mathrm{J}\) , then \(\gamma_i \in \Gamma_{C}\) . Indeed, let \(a \in \mathrm{C}\) and \(b = \gamma_i(a)\) . For \(1 \leqslant j \leqslant l\) and \(j \neq i\) ,

\[
\begin{array}{r l} \langle b, \alpha_ {j} \rangle & = \langle \overline {{\omega}} _ {i} + w _ {i} w _ {0} (a), \alpha_ {j} \rangle \\ & = \langle w _ {0} (a), w _ {i} (\alpha_ {j}) \rangle > 0 \end{array}\tag{2}
\]

since \(w_0(a) \in -\mathbf{C}'\) and \(w_i(\alpha_j)\) is negative. On the other hand,

\[
\langle b, \alpha_ {i} \rangle = 1 + \left\langle w _ {0} (a), w _ {i} (\alpha_ {i}) \right\rangle \geqslant 1 + \left\langle w _ {0} (a), \tilde {\alpha} \right\rangle > 0\tag{3}
\]

since \(w_0(a) \in -\mathbf{C}'\) , \(\tilde{\alpha} - w_i(\alpha_i)\) takes negative values on \(-\mathbf{C}'\) , and \(\langle w_0(a), \tilde{\alpha} \rangle > -1\) . Finally,

\[
\langle b, \tilde {\alpha} \rangle = n _ {i} + \left\langle w _ {0} (a), w _ {i} (\tilde {\alpha}) \right\rangle = 1 + \left\langle w _ {0} (a), w _ {i} (\tilde {\alpha}) \right\rangle <   1\tag{4}
\]

since \(w_0(a) \in -\mathbf{C}'\) and \(w_i(\tilde{\alpha})\) is a positive root. The relations (2), (3) and (4) then imply that \(b \in \mathbf{C}\) , and hence that \(\gamma_i \in \Gamma_{C}\) . It is clear that the map \(i \mapsto \gamma_i\) is injective, since \(\gamma_i(0) = \overline{\omega}_i\) . Finally, let \(\gamma \in \Gamma_{C}\) with \(\gamma \neq 1\) , and put \(\gamma = tw\) with \(t \in \mathbf{P}\) and \(w \in W\) . Then \(t \neq 1\) since \(\Gamma_{C} \cap W = \{1\}\) . On the other hand, \(t(0) = \gamma(0) \in \overline{\mathbf{C}} \cap P(\mathbf{R}^{\prime})\) and Prop. 5 implies that there exists \(i \in J\) such that \(t(0) = \overline{\omega}_i\) . Then \(\gamma_i^{-1}\gamma(0) = 0\) , hence \(\gamma = \gamma_i\) since \(\Gamma_{C} \cap W = \{1\}\) . This completes the proof.

COROLLARY. The \((\bar{\omega}_i)_{i\in J}\) form a system of representatives in \(\mathrm{P}(\mathbb{R}^{\prime})\) of the non-zero elements of \(\mathrm{P}(\mathbb{R}^{\prime}) / \mathrm{Q}(\mathbb{R}^{\prime})\) .

Indeed, if we identify \(\Gamma_{C}\) with \(\mathrm{P}(\mathbb{R}^{\prime}) / \mathrm{Q}(\mathbb{R}^{\prime})\) , the element \(\gamma_{i}\) is identified with the class of \(\overline{\omega}_{i}\) mod. \(\mathrm{Q}(\mathbb{R}^{\prime})\) .

Remarks. 1) The map \(\gamma \mapsto \gamma(0)\) is a bijection from \(\Gamma_{C}\) to \(\overline{\mathrm{C}} \cap \mathrm{P}(\mathrm{R}^{\prime})\) .

\begin{enumerate}
\def\labelenumi{\arabic{enumi})}
\setcounter{enumi}{1}
\tightlist
\item
  The group G is also the normaliser of \(W_{a}\) in the group of automorphisms of E provided only with its affine structure (cf.~Exerc. 3).
\end{enumerate}

